\setlength{\baselineskip}{20pt}
\chapter{基于模型蒸馏的特征增强算法研究（DenoiseRep++）}
\label{cha:chap4}

\section{研究问题演进与章节过渡（由 DenoiseRep 到 DenoiseRep++）}
在第\ref{cha:chap3}章基础方法已满足部署效率要求，但当去噪模块必须轻量化时，其噪声分布建模能力仍存在上限。为进一步提升复杂场景下的去噪强度与泛化性能，本章提出 DenoiseRep++，重点解决“轻量结构表达受限”与“有限去噪深度不足”两个问题，并形成从基础方法到增强方法的自然过渡。

\section{教师--学生蒸馏去噪框架}
设教师去噪器与学生去噪器的噪声预测分别为 $\hat{\epsilon}_T$、$\hat{\epsilon}_S$，蒸馏损失定义为
\begin{equation}
\mathcal{L}_{distill}=\left\|\hat{\epsilon}_T-\hat{\epsilon}_S\right\|_2^2.
\end{equation}
最终训练目标由任务损失、去噪损失和蒸馏损失联合构成：
\begin{equation}
\mathcal{L}=\mathcal{L}_{task}+\lambda_d\mathcal{L}_{denoise}+\lambda_s\mathcal{L}_{distill}.
\end{equation}
该设计在训练阶段利用高容量教师提供更稳定的噪声方向监督，推理阶段仍仅保留融合后的学生主干，不增加部署开销。

\section{有限去噪深度下的增强策略}
\subsection{多轮单层去噪}
在每层参数不增加的条件下，对单层去噪映射执行少量迭代更新，等价于在局部特征空间中延长去噪轨迹。实验中通常 1--3 轮即可获得明显收益，继续增加轮数的边际收益递减。

\subsection{多步去噪融合}
借鉴扩散蒸馏思想，将教师模型的多步去噪结果压缩到学生单步映射。训练期通过多步监督增强学生表示能力，推理期通过参数折叠保持单步执行形式，实现“训练增强、部署不增负担”。

\section{无监督与有监督学习范式}
本章方法同时支持无标签与有标签训练：无监督模式仅优化去噪损失，有监督模式将任务损失与去噪损失联合优化，从而兼顾判别性与鲁棒性。\,[需要补充具体损失定义与权重设置]

\section{本章实验设计}
\subsection{消融与模块贡献}
本章实验围绕增强模块贡献展开，重点验证蒸馏与多步融合的单独作用及协同增益。

\subsection{跨任务实验设置}
实验覆盖 ReID、分类、检测、分割与车辆重识别任务，常用指标包括 mAP/Rank-1、Top-1/Top-5、AP 系列指标及 mIoU/B-IoU。所有对比在相同骨干和近似训练预算下进行。

\section{实验结果与分析}
\subsection{蒸馏模块效果}
蒸馏模块主要提升单步噪声预测质量。在 ImageNet-1k、MSMT17、COCO、ADE20K 等任务上，加入蒸馏后指标均有稳定增长，说明教师监督可有效缓解轻量去噪层的估计偏差。

\subsection{多步融合效果}
融合模块通过多轮或多步压缩增强去噪轨迹逼近能力，在噪声敏感任务（如 ReID、检测）中提升更明显。其收益来源于对层间噪声累积的抑制和对判别边界的重整。

\subsection{跨任务泛化分析}
蒸馏与融合联合使用时通常取得最佳结果，表明二者分别作用于“单步预测精度”和“有效去噪深度”两个互补维度。
\paragraph{跨任务泛化} 在分类、检测、分割和车辆重识别任务上的一致正向增益说明，增强方法并不依赖单一任务结构，而是作用于更通用的特征噪声抑制机制。对于背景复杂或域偏移明显的数据集，性能提升更为突出。
\paragraph{可视化与统计分析} 通过 t-SNE 可视化可观察到增强后特征类簇更紧凑、类间边界更清晰。进一步统计类内方差、类间距离及 Inter/Intra 比值，可见增强后类内离散度下降、类间分离度上升，与定量指标提升一致。
\paragraph{公平性对比} 为排除“训练轮数增加导致性能提升”的干扰，在相同训练预算下对比基线与增强方法。结果显示基线在延长训练后提升有限，而 DenoiseRep++ 仍保持持续收益，说明改进主要来源于方法本身而非训练时长。

\begin{table}[!htb]
\centering
\caption{DenoiseRep++ 关键消融结果（示例）}
\label{tab:drpp-ablation}
\small
\begin{tabular}{|c|c|c|c|}
\hline
任务 & 配置 & 指标1 & 指标2 \\
\hline
ImageNet-1k & Baseline & 81.82 & 95.88 \\
ImageNet-1k & +DenoiseRep & 82.13 & 96.06 \\
ImageNet-1k & +DenoiseRep++ & 82.30 & 96.19 \\
\hline
DukeMTMC-reID & Baseline & 81.20 & 87.80 \\
DukeMTMC-reID & +DenoiseRep & 82.10 & 88.70 \\
DukeMTMC-reID & +DenoiseRep++ & 82.50 & 89.10 \\
\hline
COCO & Baseline & 42.80 & 65.10 \\
COCO & +DenoiseRep & 44.30 & 67.10 \\
COCO & +DenoiseRep++ & 44.70 & 67.40 \\
\hline
ADE20K & Baseline & 34.10 & 42.20 \\
ADE20K & +DenoiseRep & 34.80 & 42.50 \\
ADE20K & +DenoiseRep++ & 35.00 & 42.60 \\
\hline
\end{tabular}
\end{table}

\section{本章小结}
本章提出 DenoiseRep++ 并在本章内完成模块消融、跨任务泛化、可视化统计和公平性对比验证。结果表明，蒸馏与多步融合能够在不增加部署复杂度的前提下显著提升去噪能力与特征判别性。
